\chapter{Statistical Mechanics of Carrier Populations}
\label{ch:statistical-mechanics-carriers}

The preceding chapter determined the one-particle states available at each
energy.  A carrier calculation also requires the equilibrium occupation
probability of every state.  The partition function and the macroscopic
constraints determine that probability.

The purpose of this chapter is to derive the sequence
\begin{equation}
  \text{microstates}
  \longrightarrow
  \text{partition function}
  \longrightarrow
  \text{occupation}
  \longrightarrow
  n,p
  \longrightarrow
  \text{charge neutrality}
  \longrightarrow
  \mu.
  \label{eq:statistical-mechanics-carrier-sequence}
\end{equation}
This provides the statistical foundation for the intrinsic and extrinsic
semiconductors treated in the following chapters
\cite{kittel2005,szeLee2012,szeNg2007,rudan2015}.

\section{From equipartition to a probability distribution}

A familiar starting point in thermal physics is the equipartition theorem.  In
the classical regime, every independent quadratic term in the energy has mean
value
\begin{equation}
  \left\langle a x^2\right\rangle
  =\frac12\kb T.
  \label{eq:classical-equipartition-quadratic-term}
\end{equation}
A free particle has three quadratic translational contributions,
\begin{equation}
  E_{\mathrm{kin}}
  =\frac{p_x^2}{2m}
   +\frac{p_y^2}{2m}
   +\frac{p_z^2}{2m},
  \label{eq:three-dimensional-classical-kinetic-energy}
\end{equation}
so
\begin{equation}
  \langle E_{\mathrm{kin}}\rangle
  =\frac32\kb T.
  \label{eq:classical-translational-equipartition}
\end{equation}
This gives an intuitive meaning to temperature: thermal energy is distributed
among accessible degrees of freedom.

Equipartition applies to classical continuous quadratic degrees of freedom.  A
band gap is a finite quantum excitation energy, and a degree of freedom becomes
thermally frozen when its energy spacing exceeds $\kb T$.  Semiconductor carrier
populations therefore require state probabilities for one electron and mean
state occupations for many electrons.

The partition function assigns equilibrium probabilities to microscopic states
with different energies.  Its classical quadratic limit reproduces
equipartition.

\section{Microstates, macrostates, and constraints}

A \emph{microstate} specifies the complete microscopic state of the system.
For one electron in a crystal, a microstate is one particular band state,
which may be labeled by band index, wave vector, and spin:
\begin{equation}
  \alpha=(n,\bm k,s),
  \qquad E_\alpha=E_{ns}(\bm k).
  \label{eq:single-electron-band-state-label}
\end{equation}
The band calculation supplies this list of one-electron energies.  A canonical
partition function describes one electron on that spectrum; a many-electron
occupation law is introduced separately.

A \emph{macrostate} specifies macroscopic constraints such as temperature,
volume, and total charge.  Statistical mechanics assigns probabilities to the
microstates compatible with those constraints.

The permitted exchanges determine the equilibrium distribution.  An isolated
system has fixed energy and particle number.  A canonical system exchanges
energy with a reservoir at fixed particle number.  A grand-canonical system
exchanges energy and particles and is specified by $T$, $V$, and chemical
potential $\mu$.

A one-electron band and a many-electron population are different levels of
description.  The canonical partition function describes one electron
thermally distributed over the one-electron states.  Fermi--Dirac statistics
describes many electrons distributed over that same list of states.
The grand-canonical description then assigns a mean occupation to each
one-electron state, while the chemical potential is adjusted so that the total
expectation obeys the specimen's fixed charge.

\section{From reservoir multiplicity to the partition function}

Consider a small system in thermal contact with a much larger reservoir.  The
combined system is isolated with fixed energy $E_{\mathrm{tot}}$.  If the small
system occupies microstate $s$ with energy $E_s$, the reservoir has energy
\begin{equation}
  E_R=E_{\mathrm{tot}}-E_s.
  \label{eq:reservoir-energy-after-system-state}
\end{equation}
Let $\Omega_R(E_R)$ be the number of reservoir microstates available at that
energy.  Every combined microstate is equally probable, so the probability of
small-system state $s$ is proportional to the reservoir multiplicity:
\begin{equation}
  \mathbb P_s\propto
  \Omega_R(E_{\mathrm{tot}}-E_s).
  \label{eq:system-probability-from-reservoir-multiplicity}
\end{equation}

Because the reservoir is large and $E_s\ll E_{\mathrm{tot}}$, expand the
logarithm of its multiplicity:
\begin{equation}
\begin{aligned}
  \ln\Omega_R(E_{\mathrm{tot}}-E_s)
  \simeq{}&\ln\Omega_R(E_{\mathrm{tot}})\\
  &-E_s
  \left.
  \frac{\partial\ln\Omega_R}{\partial E_R}
  \right|_{E_{\mathrm{tot}}}.
\end{aligned}
  \label{eq:reservoir-log-multiplicity-expansion}
\end{equation}
Use Boltzmann's entropy relation
\begin{equation}
  S_R=\kb\ln\Omega_R
  \label{eq:boltzmann-entropy-multiplicity}
\end{equation}
and the thermodynamic definition
\begin{equation}
  \frac{1}{T}
  =\left(\frac{\partial S_R}{\partial E_R}\right)_{V,N}.
  \label{eq:temperature-entropy-derivative}
\end{equation}
Then
\begin{equation}
  \frac{\partial\ln\Omega_R}{\partial E_R}
  =\frac{1}{\kb T}
  \equiv\beta.
  \label{eq:beta-from-reservoir-multiplicity}
\end{equation}
Exponentiating \cref{eq:reservoir-log-multiplicity-expansion} gives
\begin{equation}
  \Omega_R(E_{\mathrm{tot}}-E_s)
  \propto e^{-\beta E_s}.
  \label{eq:boltzmann-weight-from-reservoir}
\end{equation}
The Boltzmann factor therefore follows from the multiplicity of the reservoir.

For a system with fixed particle number whose microstates have energies $E_s$,
normalization gives the canonical probability
\begin{equation}
  \mathbb P_s
  =\frac{e^{-\beta E_s}}{Z},
  \qquad
  \beta=\frac{1}{\kb T},
  \label{eq:canonical-state-probability}
\end{equation}
where
\begin{equation}
  \boxed{
  Z(T,V,N)=\sum_s e^{-\beta E_s}
  }
  \label{eq:canonical-partition-function}
\end{equation}
is the canonical partition function.  The normalization follows immediately:
\begin{equation}
  \sum_s\mathbb P_s
  =\frac{1}{Z}\sum_s e^{-\beta E_s}=1.
  \label{eq:canonical-probability-normalization}
\end{equation}

For any observable $A$ with value $A_s$ in microstate $s$, the equilibrium
expectation is
\begin{equation}
  \langle A\rangle
  =\sum_s A_s\mathbb P_s.
  \label{eq:canonical-observable-expectation}
\end{equation}
In particular,
\begin{equation}
\begin{aligned}
  \langle E\rangle
  &=\frac{1}{Z}\sum_sE_se^{-\beta E_s}\\
  &=-\frac{\partial\ln Z}{\partial\beta}.
\end{aligned}
  \label{eq:canonical-mean-energy}
\end{equation}
The Helmholtz free energy is
\begin{equation}
  F=-\kb T\ln Z.
  \label{eq:helmholtz-from-partition-function}
\end{equation}
The partition function normalizes the state probabilities, and its derivatives
generate equilibrium thermodynamic quantities.

\subsection{Recovering equipartition from the partition function}

For one classical quadratic coordinate with energy
\begin{equation}
  E(x)=ax^2,
  \qquad a>0,
  \label{eq:one-quadratic-coordinate-energy}
\end{equation}
the continuous canonical partition integral is
\begin{equation}
\begin{aligned}
  Z_x
  &=\int_{-\infty}^{\infty}e^{-\beta ax^2}\,\dd x\\
  &=\sqrt{\frac{\pi}{\beta a}}.
\end{aligned}
  \label{eq:gaussian-quadratic-partition-integral}
\end{equation}
Its mean energy is
\begin{equation}
\begin{aligned}
  \langle ax^2\rangle
  &=-\frac{\partial\ln Z_x}{\partial\beta}\\
  &=-\frac{\partial}{\partial\beta}
    \left[
      \frac12\ln\pi-\frac12\ln\beta-\frac12\ln a
    \right]\\
  &=\frac{1}{2\beta}
  =\frac12\kb T.
\end{aligned}
  \label{eq:equipartition-from-partition-function}
\end{equation}
Thus the partition function reproduces equipartition when the energy is
classical, continuous, and quadratic.  Unlike equipartition, the partition
function also applies to discrete and quantum spectra, which is why it is the
appropriate starting point for electronic occupations.

\subsection{Two-state canonical example}

Let a fixed particle occupy either a ground state of energy $0$ or an excited
state of energy $\Delta$.  Then
\begin{equation}
  Z=1+e^{-\beta\Delta}.
  \label{eq:two-state-canonical-partition}
\end{equation}
The excited-state probability is
\begin{equation}
  \mathbb P_{\mathrm{exc}}
  =\frac{e^{-\beta\Delta}}{1+e^{-\beta\Delta}}.
  \label{eq:two-state-excited-probability}
\end{equation}
At low temperature, $\beta\Delta\gg1$ and the excitation is rare.  At high
temperature, the two states approach equal probability.  This elementary
example already contains the central logic: enumerate microstates, assign
Boltzmann weights, sum them to form $Z$, and normalize.

\section{One electron distributed over a band}

Let $E_\alpha$ be the one-electron energies obtained from a band calculation.
For exactly one electron in thermal equilibrium, the canonical partition
function is
\begin{equation}
  Z_1=\sum_\alpha e^{-\beta E_\alpha}.
  \label{eq:single-electron-band-partition-function}
\end{equation}
If the distinct energies $E_j$ have multiplicities $M_j$, the same sum is
\begin{equation}
  Z_1=\sum_jM_je^{-\beta E_j}.
  \label{eq:single-electron-partition-with-multiplicity}
\end{equation}
The probability that the electron occupies one particular state $\alpha$ is
\begin{equation}
  \mathbb P_\alpha
  =\frac{e^{-\beta E_\alpha}}{Z_1},
  \label{eq:single-electron-canonical-state-probability}
\end{equation}
and the probability that its energy is $E_j$ is
\begin{equation}
  \mathbb P(E_j)
  =\frac{M_je^{-\beta E_j}}{Z_1}.
  \label{eq:single-electron-energy-probability-multiplicity}
\end{equation}
The factor $M_j$ arises because any one of $M_j$ distinct states gives the same
energy.  This one-electron calculation uses the canonical Boltzmann weights directly.

In a macroscopic band, the sum becomes
\begin{equation}
  Z_1
  =\mathcal V_d\int g(E)e^{-\beta E}\,\dd E,
  \label{eq:single-electron-partition-dos-integral}
\end{equation}
where $g(E)$ is the one-electron density of states per physical size.  This
equation connects the discrete multiplicities of the preceding chapter to the
canonical statistics of one electron.

\section{Many electrons and the chemical potential}

If a microstate can contain different particle numbers $N_s$, the reservoir
must account for both energy and particles.  The statistical weight becomes
\begin{equation}
  e^{-\beta(E_s-\mu N_s)}.
  \label{eq:grand-canonical-weight}
\end{equation}
The grand partition function is
\begin{equation}
  \boxed{
  \Xi(T,V,\mu)
  =\sum_s e^{-\beta(E_s-\mu N_s)}.
  }
  \label{eq:grand-partition-function}
\end{equation}
The probability of microstate $s$ is
\begin{equation}
  \mathbb P_s
  =\frac{e^{-\beta(E_s-\mu N_s)}}{\Xi}.
  \label{eq:grand-canonical-probability}
\end{equation}

The factor $-\mu N_s$ controls the statistical cost of changing particle
number.  Adding a particle with energy $E$ changes the exponent by
\begin{equation}
  -\beta(E-\mu).
  \label{eq:add-particle-grand-energy}
\end{equation}
A state far above $\mu$ is expensive to occupy, while a state far below $\mu$
is strongly favored.

The grand potential is
\begin{equation}
  \Omega=-\kb T\ln\Xi,
  \label{eq:grand-potential}
\end{equation}
and the expected particle number follows from
\begin{equation}
  \boxed{
  \langle N\rangle
  =\frac{1}{\beta}
   \frac{\partial\ln\Xi}{\partial\mu}.
  }
  \label{eq:particle-number-from-grand-partition}
\end{equation}
To verify the sign and factor, differentiate directly:
\begin{equation}
\begin{aligned}
  \frac{\partial\ln\Xi}{\partial\mu}
  &=\frac{1}{\Xi}
    \sum_s\beta N_s e^{-\beta(E_s-\mu N_s)}\\
  &=\beta\langle N\rangle.
\end{aligned}
  \label{eq:grand-particle-derivative}
\end{equation}

The fundamental thermodynamic relation for one particle species is
\begin{equation}
  \dd U=T\,\dd S-P\,\dd V+\mu\,\dd N.
  \label{eq:fundamental-relation-with-chemical-potential}
\end{equation}
Therefore
\begin{equation}
  \mu
  =\left(\frac{\partial U}{\partial N}\right)_{S,V}.
  \label{eq:chemical-potential-internal-energy-derivative}
\end{equation}
For the Helmholtz free energy $F=U-TS$,
\begin{equation}
  \dd F=-S\,\dd T-P\,\dd V+\mu\,\dd N,
  \label{eq:helmholtz-differential-chemical-potential}
\end{equation}
so the form used at fixed temperature is
\begin{equation}
  \boxed{
  \mu
  =\left(\frac{\partial F}{\partial N}\right)_{T,V}.
  }
  \label{eq:chemical-potential-free-energy-derivative}
\end{equation}

\begin{definition}{Chemical potential}{chemical-potential-statistical}
The chemical potential $\mu$ is the change in the appropriate thermodynamic
energy when the particle number changes while the natural external variables
are held fixed.  It is the variable conjugate to particle number.  In the
grand-canonical weight, $E-\mu$ is the statistical cost of placing an electron
in a one-electron state of energy $E$.
\end{definition}

The derivative is understood in the thermodynamic limit, where changing $N$ by
one is represented by a continuous derivative.  In a semiconductor, charge
neutrality or a prescribed total electron number fixes the value of $\mu$.

\begin{definition}{Fermi level}{fermi-level-definition}
At thermal equilibrium, the Fermi level is the semiconductor-physics name for
the electronic chemical potential $\mu(T)$.  For a nondegenerate spin-resolved
state governed by the Fermi--Dirac occupation,
\begin{equation}
  f(E=\mu,T)=\frac12.
  \label{eq:fermi-level-half-occupation-definition}
\end{equation}
The Fermi level can lie in a band gap, where no one-electron state exists.
\end{definition}

\begin{definition}{Fermi energy}{fermi-energy-definition}
The Fermi energy is the zero-temperature electronic chemical potential:
\begin{equation}
  \boxed{E_F=\mu(T=0).}
  \label{eq:fermi-energy-zero-temperature-definition}
\end{equation}
For a noninteracting metal with a partially filled band, it is the boundary
between occupied and empty one-electron states at zero temperature.
\end{definition}

In a gapped system, the zero-temperature particle number can remain unchanged
while $\mu$ moves within an interval containing no states.  The chemical-
potential definition remains primary; the phrase ``highest occupied energy''
is a special metallic interpretation of $E_F$.  Some semiconductor texts use
\emph{Fermi energy} and \emph{Fermi level} interchangeably at finite
temperature.  This manuscript reserves $E_F$ for the zero-temperature quantity
and uses $\mu(T)$ for the finite-temperature chemical potential.

\section{Mean occupation of one spin-resolved band state}

A semiconductor contains many electrons even when its bands are calculated
from an effective one-electron Hamiltonian.  The independent-quasiparticle
model treats the resulting band states as available one-electron orbitals and
constructs many-electron configurations by assigning occupations to them.
The effective band model incorporates its selected interaction effects into the
one-electron energies.  The independent-quasiparticle treatment then assigns
occupations while omitting residual correlations.

\begin{boundarybox}
The canonical distribution gives the state probabilities of one electron.
Fermi--Dirac statistics gives the mean occupations when many independent
fermionic quasiparticles fill the same one-electron spectrum.  The one-electron
energies and the many-electron occupation law are separate assumptions.
\end{boundarybox}

For one complete spin-resolved state of energy $E$, fermionic antisymmetry
permits two occupation possibilities:
\begin{center}
  \begin{tabular}{ccc}
    \toprule
    Occupation $n$ & Energy & Particle number \\
    \midrule
    $0$ & $0$ & $0$ \\
    $1$ & $E$ & $1$ \\
    \bottomrule
  \end{tabular}
\end{center}
Their grand-canonical weights are
\begin{equation}
  w_0=1,
  \qquad
  w_1=e^{-\beta(E-\mu)}.
  \label{eq:single-state-grand-weights}
\end{equation}
The single-state grand partition function is
\begin{equation}
  \Xi_E=1+e^{-\beta(E-\mu)}.
  \label{eq:single-state-grand-partition-statmech}
\end{equation}
The expected occupation is
\begin{equation}
\begin{aligned}
  \langle n_E\rangle
  &=0\frac{w_0}{\Xi_E}+1\frac{w_1}{\Xi_E}\\
  &=\frac{e^{-\beta(E-\mu)}}
          {1+e^{-\beta(E-\mu)}}\\
  &=\frac{1}{1+e^{\beta(E-\mu)}}.
\end{aligned}
  \label{eq:fermi-dirac-from-partition-function}
\end{equation}
Thus
\begin{equation}
  \boxed{
  f(E,\mu,T)
  =\frac{1}{1+\exp[(E-\mu)/(\kb T)]}.
  }
  \label{eq:fermi-dirac-statistical-derivation}
\end{equation}
The plus sign in the denominator follows from the two mutually exclusive
occupations $n=0$ and $n=1$.  This restriction belongs to the many-electron filling of the one-electron
spectrum.  The band $E_n(\bm k)$ and its density of states are calculated at the
one-electron level.

The probability that the state is empty is
\begin{equation}
\begin{aligned}
  1-f(E,\mu,T)
  &=\frac{1}{1+\exp[(\mu-E)/(\kb T)]}.
\end{aligned}
  \label{eq:empty-fermion-state-probability-statmech}
\end{equation}
This vacancy probability will become the hole occupation in a valence band.

Three limits provide immediate checks:
\begin{align}
  E=\mu&:\quad f=\frac12,
  \label{eq:statmech-half-occupation}\\
  E-\mu\gg\kb T&:\quad f\rightarrow0,
  \label{eq:statmech-empty-high-state}\\
  \mu-E\gg\kb T&:\quad f\rightarrow1.
  \label{eq:statmech-filled-low-state}
\end{align}
The statement $f=1/2$ at $E=\mu$ applies to one nondegenerate band state.  A
localized donor or acceptor with unequal charge-state degeneracies acquires an
additional statistical factor, as shown later.

\section{Many independent fermionic states}

Let the one-particle energies be $E_i$.  A many-electron microstate is specified
by occupations $n_i\in\{0,1\}$, with
\begin{equation}
  E_s=\sum_i n_iE_i,
  \qquad
  N_s=\sum_i n_i.
  \label{eq:independent-fermion-microstate-energy-number}
\end{equation}
The grand partition function is
\begin{equation}
\begin{aligned}
  \Xi
  &=\sum_{\{n_i\}}
    \exp\left[-\beta\sum_i n_i(E_i-\mu)\right]\\
  &=\sum_{\{n_i\}}
    \prod_i e^{-\beta n_i(E_i-\mu)}.
\end{aligned}
  \label{eq:many-fermion-grand-sum}
\end{equation}
Because each $n_i$ can be summed independently,
\begin{equation}
\begin{aligned}
  \Xi
  &=\prod_i\sum_{n_i=0}^{1}
    e^{-\beta n_i(E_i-\mu)}\\
  &=\boxed{
    \prod_i\left[1+e^{-\beta(E_i-\mu)}\right]}.
\end{aligned}
  \label{eq:ideal-fermion-grand-partition-product}
\end{equation}
Taking the logarithm turns the product into a sum:
\begin{equation}
  \ln\Xi
  =\sum_i\ln\left[1+e^{-\beta(E_i-\mu)}\right].
  \label{eq:log-grand-partition-fermions}
\end{equation}
Using \cref{eq:particle-number-from-grand-partition},
\begin{equation}
\begin{aligned}
  \langle N\rangle
  &=\frac{1}{\beta}
    \frac{\partial}{\partial\mu}\ln\Xi\\
  &=\sum_i\frac{1}{1+e^{\beta(E_i-\mu)}}\\
  &=\sum_i f(E_i,\mu,T).
\end{aligned}
  \label{eq:particle-number-sum-of-fermi-occupations}
\end{equation}
The total population is therefore a sum of single-state occupation
probabilities.

\section{From a state sum to carrier concentrations}

The density of states converts the discrete state sum into an energy integral:
\begin{equation}
  \sum_i A(E_i)
  \longrightarrow
  \mathcal V_d\int g(E)A(E)\,\dd E.
  \label{eq:state-sum-to-dos-integral}
\end{equation}
For conduction-band electrons in a bulk semiconductor,
\begin{equation}
  \boxed{
  n(\mu,T)
  =\int_{E_C}^{\infty}
   g_c(E)f(E,\mu,T)\,\dd E.
  }
  \label{eq:statmech-electron-concentration}
\end{equation}
This equation has two independent ingredients:
\begin{equation}
  \underbrace{g_c(E)}_{\text{available states}}
  \times
  \underbrace{f(E,\mu,T)}_{\text{occupied fraction}}.
  \label{eq:states-times-occupation}
\end{equation}
Their product is the occupied-state density per energy.

A hole is a missing electron in an otherwise occupied valence band.  The
probability that a valence state is missing its electron is $1-f$, so
\begin{equation}
  \boxed{
  p(\mu,T)
  =\int_{-\infty}^{E_V}
   g_v(E)[1-f(E,\mu,T)]\,\dd E.
  }
  \label{eq:statmech-hole-concentration}
\end{equation}
The hole expression counts vacancies in electronic valence states.

\begin{modelcheck}
If $g(E)$ is in states per electronvolt per cubic centimetre and $\dd E$ is in
electronvolts, then the integral gives carriers per cubic centimetre.  If the
DOS is instead per joule per cubic metre, both the energy and spatial units must
change together.  The occupation probability is dimensionless.
\end{modelcheck}

\section{The Boltzmann limit derived from the partition function}

For a state several thermal energies above the chemical potential, define
\begin{equation}
  y=\frac{E-\mu}{\kb T}\gg1.
  \label{eq:statmech-boltzmann-reduced-energy}
\end{equation}
Then
\begin{equation}
\begin{aligned}
  f(E)
  &=\frac{e^{-y}}{1+e^{-y}}\\
  &\simeq e^{-y}
  =\exp\left[-\frac{E-\mu}{\kb T}\right].
\end{aligned}
  \label{eq:statmech-boltzmann-tail}
\end{equation}
This is the low-occupation limit of the fermionic partition function.  The
relative overestimate made by using
$e^{-y}$ is
\begin{equation}
  \frac{e^{-y}-f}{f}=e^{-y}.
  \label{eq:statmech-boltzmann-relative-error}
\end{equation}
At $y=3$ it is about five percent; at $y=5$ it is below one percent.

Using a three-dimensional parabolic conduction-band DOS,
\begin{equation}
  g_c(E)=A_c\sqrt{E-E_C},
  \label{eq:statmech-parabolic-conduction-dos}
\end{equation}
and the Boltzmann occupation gives
\begin{equation}
\begin{aligned}
  n
  &\simeq\int_{E_C}^{\infty}
    A_c\sqrt{E-E_C}
    e^{-(E-\mu)/(\kb T)}\,\dd E\\
  &=e^{-(E_C-\mu)/(\kb T)}
    A_c\int_0^\infty
    \sqrt{\varepsilon}
    e^{-\varepsilon/(\kb T)}\,\dd\varepsilon.
\end{aligned}
  \label{eq:statmech-boltzmann-electron-integral}
\end{equation}
Let $x=\varepsilon/(\kb T)$.  Then
\begin{equation}
\begin{aligned}
  \int_0^\infty\sqrt{\varepsilon}
  e^{-\varepsilon/(\kb T)}\,\dd\varepsilon
  &=(\kb T)^{3/2}
    \int_0^\infty x^{1/2}e^{-x}\,\dd x\\
  &=(\kb T)^{3/2}\Gamma\left(\frac32\right)\\
  &=(\kb T)^{3/2}\frac{\sqrt\pi}{2}.
\end{aligned}
  \label{eq:statmech-gamma-integral}
\end{equation}
Collecting the state-counting factors defines the effective conduction density
of states $N_C$ and yields
\begin{equation}
  \boxed{
  n=N_C\exp\left[-\frac{E_C-\mu}{\kb T}\right].
  }
  \label{eq:statmech-nondegenerate-electron-result}
\end{equation}
The exponential comes from the partition function; $N_C$ comes from integrating
the band density of states.  Confusing these two origins makes the formula seem
arbitrary.

The same reasoning for valence-band vacancies gives
\begin{equation}
  \boxed{
  p=N_V\exp\left[-\frac{\mu-E_V}{\kb T}\right].
  }
  \label{eq:statmech-nondegenerate-hole-result}
\end{equation}

\section{The chemical potential is selected by a constraint}

Equations~\eqref{eq:statmech-electron-concentration} and
\eqref{eq:statmech-hole-concentration} give populations for a trial $\mu$.
A particle-number or charge constraint closes these population equations and
determines $\mu$.

For an intrinsic neutral semiconductor,
\begin{equation}
  n(\mu,T)=p(\mu,T).
  \label{eq:statmech-intrinsic-neutrality}
\end{equation}
For a uniform semiconductor with ionized donors and acceptors,
\begin{equation}
  n(\mu,T)+N_A^-(\mu,T)
  =p(\mu,T)+N_D^+(\mu,T).
  \label{eq:statmech-extrinsic-neutrality}
\end{equation}
The equilibrium chemical potential is the value that satisfies the appropriate
constraint.  Thus a ``Fermi-level calculation'' is fundamentally a constrained
statistical-mechanics calculation.

Under nondegenerate conditions, multiplication of
\cref{eq:statmech-nondegenerate-electron-result,eq:statmech-nondegenerate-hole-result}
gives
\begin{equation}
  np
  =N_CN_Ve^{-(E_C-E_V)/(\kb T)}.
  \label{eq:statmech-carrier-product}
\end{equation}
Defining
\begin{equation}
  n_i^2=N_CN_Ve^{-E_g/(\kb T)}
  \label{eq:statmech-intrinsic-product-definition}
\end{equation}
produces
\begin{equation}
  np=n_i^2.
  \label{eq:statmech-mass-action-law}
\end{equation}
The mass-action law is therefore a consequence of equilibrium, parabolic-band
state counting, and the Boltzmann limit.  Degenerate populations require the
full Fermi--Dirac carrier integrals.

\section{Localized states and internal degeneracy}

A localized dopant has a finite set of charge configurations.  The partition
sum includes the internal degeneracy of each configuration.

Suppose a donor has an ionized configuration $D^+$ with reference weight one
and a neutral configuration $D^0$ with degeneracy ratio $g_D$.  Their weights
are
\begin{equation}
  w(D^+)=1,
  \qquad
  w(D^0)=g_D e^{(\mu-E_D)/(\kb T)}.
  \label{eq:statmech-donor-configuration-weights}
\end{equation}
Normalization gives
\begin{equation}
  \mathbb P(D^+)
  =\frac{1}{1+g_De^{(\mu-E_D)/(\kb T)}}.
  \label{eq:statmech-donor-ionized-probability}
\end{equation}
Therefore
\begin{equation}
  N_D^+=N_D\mathbb P(D^+).
  \label{eq:statmech-ionized-donor-concentration}
\end{equation}
The degeneracy factor originates in the partition sum over microscopic
configurations.  Acceptor
statistics follow from the same procedure with their own charge states and
degeneracies.  The extrinsic-semiconductor chapter develops both cases fully.

\section{A disciplined route through a carrier calculation}

A transparent equilibrium calculation begins with the bands
$E_n(\bm k)$ and their degeneracies.  Their spectral multiplicities determine
the densities of available states.  Partition sums for band and localized
states determine occupation and vacancy probabilities.  Energy integration
then gives $n$, $p$, and the charged dopant populations.  Total particle number
or charge neutrality fixes $\mu$.  Normalization, limiting behavior, and the
residual of the closing constraint verify the result.

\section{Chapter summary}

The canonical partition function normalizes Boltzmann weights for fixed particle
number.  For one electron, it assigns probabilities directly to the
one-electron band states.  Particle exchange in a many-electron subsystem
introduces the grand partition
function and the chemical potential.  Summing the empty and occupied
configurations of each spin-resolved state produces the Fermi--Dirac mean
occupation.  For independent states,
the grand partition function factorizes and the expected particle number is the
sum of those occupations.  The density of states converts that sum into the
carrier integrals for conduction electrons and valence-band holes.  The
Boltzmann carrier formulas follow from the low-occupation limit and the
integrated parabolic density of states.  Finally, charge neutrality selects
the chemical potential.  This chain is the statistical-mechanical foundation of
all equilibrium carrier-concentration formulas used in the following chapters.
